\FloatBarrier
\newpage
\section{导数}
\label{sec:calculus-derivatives}

\begingroup
\setlength{\parskip}{5pt}

\subsection{平均变化率}

设自变量由\(x_0\)变到\(x_0+h\)。相应的自变量增量与函数增量分别为
\begin{equation*}
 \Delta x=h,\qquad \Delta y=f(x_0+h)-f(x_0).
\end{equation*}
当\(h\ne0\)时，称
\begin{equation}\label{eq:calculus-average-rate}
 \frac{\Delta y}{\Delta x}=\frac{f(x_0+h)-f(x_0)}h
\end{equation}
为这一段上的{\heiti{}平均变化率}。在函数图像上，它是两点\((x_0,f(x_0))\)、\((x_0+h,f(x_0+h))\)所确定割线的斜率；当\(f\)表示位置、\(x\)表示时间时，它就是平均速度。

例如，\(f(x)=x^2\)时，
\begin{equation*}
 \frac{(x_0+h)^2-x_0^2}{h}=2x_0+h.
\end{equation*}
固定\(x_0\)，缩短所考察的范围，平均变化率就趋于\(2x_0\)；改变\(x_0\)，这个数也随之改变。前者是在一个位置附近确定变化率，后者则是比较不同位置处的变化。

\begin{figure}[!ht]
\centering
\includegraphics[width=.43\linewidth]{calculus-secants}
\caption{固定点\(P\)，另一点沿曲线趋近它；割线斜率趋于切线\(\ell\)的斜率。}
\label{fig:calculus-secants}
\end{figure}

\begin{example}
某质点的位置为\(s(t)=t^2+2t\)。分别计算从\(t=1\)到\(t=1+h\)及从\(t=2\)到\(t=2+h\)的平均速度。令\(h\)分别取正数和负数，解释所比较的时间范围。
\end{example}

\subsection{一点处的导数}

\begin{definition}
设\(f\)在\(x_0\)附近有定义。若有限极限
\begin{equation}\label{eq:calculus-derivative-definition}
 f'(x_0)=\lim_{h\to0}\frac{f(x_0+h)-f(x_0)}h
\end{equation}
存在，则称\(f\)在\(x_0\)处{\heiti{}可导}，这个极限称为\(f\)在该点的{\heiti{}导数}。
\end{definition}

令\(x=x_0+h\)，定义也可以写为
\begin{equation*}
 f'(x_0)=\lim_{x\to x_0}\frac{f(x)-f(x_0)}{x-x_0}.
\end{equation*}
这里始终固定\(x_0\)，让另一个点向它靠近。瞬时速度由这样的极限给出；它的确定并不需要令时间增量等于零。

只取\(h\to0^-\)或\(h\to0^+\)，分别得到左导数与右导数。函数在内点可导，当且仅当两者都存在且相等。区间端点处需要导数时，按相应的单侧导数理解。

由定义直接得到，常值函数的导数为零，\(f(x)=kx+b\)在每一点的导数均为\(k\)，而\(f(x)=x^2\)在\(x_0\)处的导数为\(2x_0\)。对于\(f(x)=|x|\)，原点处的差商为\(|h|/h\)，左右极限分别为\(-1\)与\(1\)，所以它在原点不可导。

\begin{example}
用定义求\(f(x)=x^3\)在\(x_0\)处的导数。再设\(g(x)=\sqrt{x}\)，用有理化求\(g'(x_0)\)，其中\(x_0>0\)。
\end{example}

\begin{exercise}
\mbox{}\par\nobreak

（1）用定义求\(f(x)=1/x\)在\(x_0\ne0\)处的导数。

（2）设\(f(x)=x|x|\)。判断它在原点是否可导，并与\(|x|\)比较。

（3）设\(f(0)=0\)，\(x\ne0\)时\(f(x)=x\sin(1/x)\)。判断\(f\)在原点是否连续、可导。
\end{exercise}

\subsection{可导函数的连续性}

可导性研究的是函数增量与自变量增量的比。若这个比趋于有限的数，那么自变量增量趋零时，函数增量本身也必趋零。

\begin{theorem}\label{thm:calculus-differentiable-continuous}
若\(f\)在\(x_0\)处可导，则它在该点连续。
\end{theorem}

\begin{proof}
当\(h\ne0\)时，
\begin{equation*}
 f(x_0+h)-f(x_0)=\frac{f(x_0+h)-f(x_0)}h\,h\longrightarrow f'(x_0)\cdot0=0.
\end{equation*}
因此\(f(x_0+h)\to f(x_0)\)。
\end{proof}

反过来，连续只要求函数增量趋零，并未规定它与\(h\)之间的比例。\(|x|\)在原点连续而不可导，就是一个例子。

比较\(|x|\)、\(x\sin(1/x)\)与\(x^2\sin(1/x)\)，后两者在原点补定义为零。三者都在原点连续，但差商分别为\(|h|/h\)、\(\sin(1/h)\)、\(h\sin(1/h)\)，只有最后一个趋于零。

\subsection{导函数}

\begin{definition}
将\(f\)在各可导点处的导数与该点对应，得到的函数称为\(f\)的{\heiti{}导函数}，记为\(f'\)。其定义域是\(f\)所有可导点组成的集合。
\end{definition}

例如，\(f(x)=x^2\)的导函数为\(f'(x)=2x\)。原函数的图像是一条抛物线，导函数的图像则是一条直线；后者的纵坐标记录前者在相应横坐标处的变化率。\(f'(1)=2\)是一个数，\(f'\)则是一个新的函数。

若记\(y=f(x)\)，也常写
\begin{equation*}
 y'=f'(x)=\frac{dy}{dx},\qquad \frac{d}{dx}f(x)=f'(x).
\end{equation*}
其中\(d/dx\)表示对\(x\)求导。正的导数与函数上升、负的导数与函数下降之间的联系，可以从图像观察；要把这一观察推进为区间上的判据，还需要后面的中值定理。

\begin{example}
求\(f(x)=|x|\)的导函数及其定义域。已知某函数在\((-1,0)\)上导数为负、在\((0,1)\)上导数为正，试画出可能的原函数图像，并考虑\(x=0\)处的取值是否已被这些信息确定。
\end{example}

\subsection{切线}

图像上经过\(P=(x_0,f(x_0))\)与\(Q=(x_0+h,f(x_0+h))\)的割线，可以取方向向量
\begin{equation*}
 \left(1,\frac{f(x_0+h)-f(x_0)}h\right).
\end{equation*}
若\(f\)在\(x_0\)处可导，割线方向趋于\((1,f'(x_0))\)，所以切线方程为
\begin{equation}\label{eq:calculus-tangent}
 y-f(x_0)=f'(x_0)(x-x_0).
\end{equation}
法线的方向向量可取\((-f'(x_0),1)\)，因而其方程可写为
\begin{equation*}
 (x-x_0)+f'(x_0)(y-f(x_0))=0.
\end{equation*}
切线由局部方向确定，不要求图像始终位于直线同一侧。例如，\(y=x^3\)在原点的切线为\(y=0\)，图像却穿过这条切线。

\begin{example}
用导数求抛物线\(y=x^2\)在\((1,1)\)处的切线和法线，并用第四章的方法核对。再求\(y=x^3\)在原点的切线。
\end{example}

\begin{example}[2012年广东卷理科第12题]
\label{ex:calculus-gaokao-polynomial-tangent}
求曲线\(y=x^3-x+3\)在点\((1,3)\)处的切线方程。用导数定义确定切线斜率。
\end{example}

\begin{exercise}
\mbox{}\par\nobreak

（1）（2013年大纲卷文科第10题改编）曲线\(y=x^4+ax^2+1\)在点\((-1,a+2)\)处切线的斜率为\(8\)，求\(a\)。用差商极限完成计算。

（2）（2012年安徽卷文科第17题第（2）问）设\(f(x)=ax+1/(ax)+b\)，\(a>0\)。若曲线\(y=f(x)\)在点\((1,f(1))\)处的切线方程为\(y=3x/2\)，求\(a,b\)。
\end{exercise}

\subsection{微分}
\label{subsec:calculus-differential}

平方函数的增量满足
\begin{equation*}
 (x_0+h)^2-x_0^2=2x_0h+h^2.
\end{equation*}
其中\(2x_0h\)关于\(h\)线性，而余下的\(h^2\)与\(h\)之比趋零。一般可导函数也有这种结构。

\begin{proposition}\label{prop:calculus-first-order}
函数\(f\)在\(x_0\)处可导，当且仅当存在常数\(a\)，使
\begin{equation}\label{eq:calculus-first-order}
 f(x_0+h)=f(x_0)+ah+r(h),\qquad \lim_{h\to0}\frac{r(h)}h=0.
\end{equation}
此时\(a=f'(x_0)\)。
\end{proposition}

\begin{proof}
若可导，取\(a=f'(x_0)\)，将定义中的差商减去\(a\)，就得到\(r(h)/h\to0\)。反过来，将式\eqref{eq:calculus-first-order}减去\(f(x_0)\)并除以\(h\)，差商便趋于\(a\)。
\end{proof}

\begin{definition}
对可导函数\(y=f(x)\)，称\(dx=\Delta x\)为自变量的微分，称函数增量的线性部分
\begin{equation*}
 dy=f'(x)\,dx
\end{equation*}
为函数在\(x\)处的{\heiti{}微分}。
\end{definition}

因此，一次近似为\(f(x+h)\approx f(x)+f'(x)h\)。\(\Delta y\)是实际增量，\(dy\)是线性近似，两者之差与\(\Delta x\)之比趋零。导数记号\(dy/dx\)也就与这个线性关系相容。

\begin{figure}[!ht]
\centering
\includegraphics[width=.47\linewidth]{calculus-differential}
\caption{同一自变量增量下，\(\Delta y\)由曲线确定，\(dy\)由切线确定，两者之差为余项。}
\label{fig:calculus-differential}
\end{figure}

例如，\(x=4\)时，\(\sqrt{x}\)的导数为\(1/4\)，于是
\begin{equation*}
 \sqrt{4+h}=2+\frac h4+r(h),\qquad \frac{r(h)}h\to0.
\end{equation*}
还可以由有理化直接检验误差：
\begin{equation*}
 \sqrt{4+h}-2-\frac h4=-\frac{h^2}{4(\sqrt{4+h}+2)^2}.
\end{equation*}
近似值的意义与误差的控制，在这里来自同一个增量表达式。

一次近似也可以估计测量误差。正方形边长由\(a\)变为\(a+h\)，面积的实际增量为\(2ah+h^2\)，微分为\(2ah\)。当\(a>0\)固定、\(|h|\)很小时，面积的相对变化约为\(2h/a\)，是边长相对变化的两倍。

若\(f'(x_0)=0\)，一次近似成为\(f(x_0+h)\approx f(x_0)\)。这表示增量比\(h\)更快趋零，并不表示函数在附近恒定。\(x^2\)与\(x^3\)在原点的线性部分都为零，实际增量却分别为\(h^2\)、\(h^3\)。要继续区分这些变化，后面需要更高阶的近似。

\begin{exercise}
\mbox{}\par\nobreak

（1）用微分近似计算\(\sqrt{4.04}\)，再用上式给出误差上界。

（2）对\(f(x)=1/x\)，写出在\(x=1\)处的增量、微分及两者之差。

（3）设\(f(0)=0\)，\(x\ne0\)时\(f(x)=x^2\sin(1/x)\)。证明\(f\)在原点可导，并写出其在原点的一次近似。

（4）设\(f(a)=2\)、\(f'(a)=3\)。分别写出\(f(a+2h)\)与\(f(a-h)\)在\(h\to0\)时的一次近似，并求\(\displaystyle\lim_{h\to0}\frac{f(a+2h)-f(a-h)}h\)。

（5）圆的半径为\(r>0\)，测量误差为\(h\)。写出面积的实际增量与微分；若半径的相对误差不超过\(1\%\)，估计面积的相对误差，并给出由精确增量得到的上界。

（6）设\(f\)在\(a\)处可导，证明\(\bigl(f(a+h)+f(a-h)-2f(a)\bigr)/h\to0\)。这一结论能否推出\(f\)在\(a\)附近为常值？

（7）设\(x<0\)时\(f(x)=ax+b\)，\(x\ge0\)时\(f(x)=x^2\)。分别求\(f\)在原点连续、可导的条件，并用增量的线性部分解释所得结果。
\end{exercise}

\endgroup
\FloatBarrier
\newpage
\subsection{向量值函数的微分}\label{subsec:calculus-vector-differential}

回到第\ref{subsec:vector-valued-functions}小节的动点。若位置由 \(\vec r(t)\) 给出，从时刻 \(t\) 到 \(t+h\) 的位移就是 \(\vec r(t+h)-\vec r(t)\)。取第三章的抛物线 \(\vec r(t)=(t,t^2)\)，直接相减得
\[
 \vec r(t+h)-\vec r(t)=(h,2th+h^2)=h(1,2t)+(0,h^2).
\]
先看速度。将位移除以 \(h\)，得到平均速度 \((1,2t+h)\)；令 \(h\to0\)，就得到时刻 \(t\) 的瞬时速度 \((1,2t)\)。其中 \(2t\) 只表示纵坐标的变化率；速度是整个向量，它的模 \(\sqrt{1+4t^2}\) 是速率。

一般地，设 \(\vec r(t)=(x(t),y(t))\)。若下面的向量极限存在，就称它为 \(\vec r\) 在 \(t\) 处的{\heiti{}导数}：
\[
 \vec r\,'(t)=\lim_{h\to0}\frac{\vec r(t+h)-\vec r(t)}h.
\]
差商的两个分量正是 \(x,y\) 各自的差商。由第三章的向量极限判据，\(\vec r\) 可导当且仅当两个分量都可导，此时 \(\vec r\,'(t)=(x'(t),y'(t))\)。当导向量非零时，它给出轨迹在该处的切向。

再看位移。上例中 \(h(1,2t)\) 随时间增量 \(h\) 线性变化，余下的 \((0,h^2)\) 的模除以 \(|h|\) 趋于零。一般地，从导数的定义即得
\[
 \vec r(t+h)-\vec r(t)=h\vec r\,'(t)+\vec R(h),\qquad
 \frac{|\vec R(h)|}{|h|}\longrightarrow0.
\]
取 \(dt=h\)，称位移的线性部分 \(d\vec r=\vec r\,'(t)\,dt\) 为{\heiti{}微分}。在动点的情形，它就是按当前速度运动 \(h\) 个时间单位所产生的位移。

第三章的例\ref{ex:vector-position-segment}中，两个点沿同一条线段从 \(A\) 走向 \(B\)。记 \(\vec u=\overrightarrow{AB}\)，它们的位置分别为 \(\vec r(t)=\vec a+t\vec u\)、\(\vec s(t)=\vec a+t^2\vec u\)，其中 \(0\leq t\leq1\)。直接比较一段时间内的位移：
\[
 \vec r(t+h)-\vec r(t)=h\vec u,\qquad
 \vec s(t+h)-\vec s(t)=2th\vec u+h^2\vec u.
\]
第一个点的速度恒为 \(\vec u\)，微分恰好等于实际位移；第二个点的速度为 \(2t\vec u\)，用微分估计位移时，相差 \(h^2\vec u\)。两点先后经过中点时，前者的速度是 \(\vec u\)，后者的速度是 \(\sqrt2\,\vec u\)：同一位置、同一方向，速率却不同。

\begin{exercise}
\mbox{}\par\nobreak

（1）对 \(\vec r(t)=(t,t^2)\)，求 \(t_0=1\) 处的位置向量的一次近似，直接检验余项条件。

（2）对参数曲线 \(\vec r(t)=(t^2,t^3)\)，分别求 \(t_0=1\) 与 \(t_0=0\) 处位移的线性部分。零点处的导向量能否给出切线方向？再由割线方向判断该处的切线。
\end{exercise}

\subsection{局部线性模型*}
\label{subsec:calculus-linear-model}

前两小节中，自变量的增量都是实数 \(h\)，函数的增量却可以是实数，也可以是向量。为什么两种微分都有同样的形式？设 \(L\) 是关于 \(h\) 的线性映射，由数乘的运算律，\(L(h)=hL(1)\)。所以只要知道单位增量的像 \(L(1)\)，整个线性映射便已确定；它的值是数还是向量，都不影响这个结论。

可导性要求的正是这样一个映射 \(L\)：用 \(L(h)\) 近似实际增量，误差的大小与 \(|h|\) 之比须趋于零。将实际增量除以 \(h\)，所得差商便趋于 \(L(1)\)，因而 \(L(1)\) 就是导数，线性主部也由此唯一确定。在基点 \(t_0\)，实值函数的微分把 \(h\) 送到 \(hf'(t_0)\)，位置函数的微分则把同一个 \(h\) 送到 \(h\vec r\,'(t_0)\)。

切线对应的映射为
\begin{equation*}
 x\longmapsto f(x_0)+L(x-x_0).
\end{equation*}
它在基点和平移后的增量之间建立联系，具有仿射形式。线性主部描述变化，基点决定这段近似放在什么位置。

例如，将函数的图像写成 \(\vec r(t)=(t,f(t))\)，则向量的一次近似给出
\[
 (x,y)=(t_0,f(t_0))+h(1,f'(t_0)).
\]
由第一分量可知 \(h=x-t_0\)，代入第二分量，恰得前面的切线方程。斜率与切向量在这里描述的是同一条直线。

\begin{exercise}
\mbox{}\par\nobreak

（1）求\(f(x)=x^3\)在\(x_0=1\)处的线性主部，并区分它与切线所给的仿射映射。

（2）回看例\ref{ex:calculus-gaokao-polynomial-tangent}，写出函数在\(x=1\)处的线性主部，并求原函数与切线在\(x=1+h\)处的高度差。当\(|h|\le1/10\)时，给出这个误差关于\(h^2\)的上界。
\end{exercise}
